\section{Methods}
\label{sec:methods}

\subsection{Construction of the Correlation Matrix}
\label{subsec:matrix-computation}
Given a molecular dynamics trajectory of length $T_{\text{traj}}$ from which we sample $M$ velocity snapshots at uniform intervals $\delta t$, we compute the finite-sample estimator of Equation~\ref{eqn:Cij} as
\begin{align}
    \hat{C}_{ij} = \delta t \sum\limits_{\tau = 0}^{\tau_{\max}/\delta t} \frac{1}{M - \tau} \sum\limits_{m = 0}^{M - \tau - 1} \mathbf{v}_{i}(t_{m} + \tau\, \delta t) \cdot \mathbf{v}_{j}(t_{m}),
    \label{eqn:Chat}
\end{align}
where the outer sum is a Riemann approximation to the time integral, truncated at a lag $\tau_{\max}$ that we choose to be several times the integral relaxation time of the velocity autocorrelation (see below), and the inner sum averages over time origins $t_{m}$.

Evaluating Equation~\ref{eqn:Chat} by direct summation costs $\mathcal{O}(|\N|^{2} M^{2})$ and is intractable at the system sizes we use.
Instead, for each pair $(i,j)$ we compute the cross-correlation in frequency space via the Wiener--Khinchin theorem, using a single pair of FFTs per Cartesian component and summing the three component contributions, reducing the cost to $\mathcal{O}(|\N|^{2} M \log M)$.
The implementation is memory-bound at large $|\N|$, so we process blocks of pair indices in sequence and reuse the forward transforms of each atom's velocity time series.

Two practical points deserve attention.
First, when $a_{i}$ is a vector such as velocity, we compute $\hat{C}_{ij}$ per Cartesian component and average; this reduces variance by a factor of three relative to using a scalar projection and preserves rotational invariance.
Second, the choice of $\tau_{\max}$ is not innocuous.
Too short a cutoff under-integrates long-lived correlations and biases $\hat{C}$ downward; too long a cutoff accumulates noise from the flat tail of the autocorrelation function and inflates the variance of every off-diagonal entry.
We estimate the integral relaxation time
\begin{align}
    \tau_{\text{int}} = \int\limits_{0}^{\infty} dt \, \frac{\langle \mathbf{v}_{i}(t) \cdot \mathbf{v}_{i}(0) \rangle}{\langle \mathbf{v}_{i}(0) \cdot \mathbf{v}_{i}(0) \rangle}
    \label{eqn:tau-int}
\end{align}
directly from the trajectory, averaging over atoms, and set $\tau_{\max} = 10\, \tau_{\text{int}}$ by default.
The same estimate provides the effective number of independent time samples $T_{\text{eff}} = T_{\text{traj}} / \tau_{\text{int}}$ used in the noise model of Section~\ref{subsec:de-noising-correlation-matrices}.

\subsection{De-noising via the Marchenko--Pastur Null}
\label{subsec:de-noising-correlation-matrices}
For a system of $|\N|$ non-interacting particles sampled at $T_{\text{eff}}$ independent times, the bulk of the eigenvalue distribution of the trace-normalised correlation matrix $\hat{C} / \mathrm{Tr}(\hat{C})$ follows the Marchenko--Pastur density~\cite{TODO:marchenko1967},
\begin{align}
    \rho_{\text{MP}}(\mu) = \frac{1}{2 \pi q\, \sigma^{2}\, \mu} \sqrt{(\mu_{+} - \mu)(\mu - \mu_{-})}, \qquad \mu \in [\mu_{-}, \mu_{+}],
    \label{eqn:mp-density}
\end{align}
with aspect ratio $q = |\N| / T_{\text{eff}}$, edges $\mu_{\pm} = \sigma^{2}\, (1 \pm \sqrt{q})^{2}$, and noise scale $\sigma^{2}$ fitted from the bulk.
We refer to $\rho_{\text{MP}}$ as the noise-only null and use it in two ways.

First, we test how close the empirical spectrum $\rho_{\text{emp}}(\mu)$ is to the null using the Kolmogorov--Smirnov distance $D_{\mathrm{KS}}(\rho_{\text{emp}}, \rho_{\text{MP}})$ and a comparison between the empirical largest eigenvalue $\mu_{\max}$ and the MP edge $\mu_{+}$.
For a system with genuine correlation structure, $D_{\mathrm{KS}}$ remains non-trivial even in the $T_{\text{eff}} \to \infty$ limit and $\mu_{\max}$ sits above $\mu_{+}$; for an uncorrelated reference, both quantities approach zero.

Second, we use the null to separate signal from noise in the entropy observable.
Eigenvalues above a threshold $\mu_{\text{cut}} = \mu_{+} (1 + \kappa / \sqrt{T_{\text{eff}}})$, with $\kappa$ a small safety factor set to $\kappa = 3$ throughout, are classified as outliers; the remainder of the spectrum is taken to be consistent with the null.
We report two entropy quantities: the raw $S(\hat{C})$ of Equation~\ref{eqn:vnentropy} and the de-noised version
\begin{align}
    S_{\text{signal}}(\hat{C}) = S(\hat{C}) - S_{\text{bulk}}(\hat{C}),
    \label{eqn:signal-entropy}
\end{align}
where $S_{\text{bulk}}$ is the von Neumann entropy computed from only those eigenvalues below $\mu_{\text{cut}}$.
$S_{\text{signal}}$ is the component of the spectrum that cannot be explained by the noise-only null and is the quantity most sensitive to the physical features we want to expose.

We validate this machinery on synthetic trajectories with prescribed correlation structure before applying it to molecular dynamics output; the validation is reported in Experiment~\ref{subsec:convergence-of-the-noise-term}.

\subsection{Molecular Dynamics Simulations}
\label{subsec:molecular-dynamics-simulations}
Simulations fall into three families, one per experimental system, and use the classical molecular dynamics code LAMMPS~\cite{TODO:lammps} unless otherwise noted.
All production runs use Langevin dynamics in the constant-$NVT$ ensemble unless $NPT$ is explicitly required, with a damping time of $0.1$~ps.
For $NPT$ runs we use a Nos\'e--Hoover barostat with coupling constants of $0.1$~ps (thermostat) and $1.0$~ps (barostat).
Each run is initialised from a randomly seeded velocity distribution and separately reported in output metadata, and every reported observable is averaged over a minimum of three and up to five independent seeds.

\paragraph{Lennard-Jones fluids (Experiments~\ref{subsec:convergence-of-the-noise-term}, \ref{subsec:bond-addition}, and~\ref{subsec:lennard-jones-phases}).}
The Lennard-Jones potential is taken with argon parameters $\sigma = 3.4$~\AA{}, $\varepsilon / k_{B} = 120$~K, and particle mass $40$~amu, cut and shifted at $r_{c} = 2.5\,\sigma$.
System size is $|\N| = 512$ for the phase sweep and the bond-addition experiment, and we sweep $|\N| \in \{64, 256, 1024, 4096\}$ in the noise-convergence study.
The integration timestep is $\delta t = 2$~fs.
For the phase sweep we equilibrate for $10^{5}$ steps in $NPT$ at $P = 1$~atm and run for $5 \times 10^{5}$ steps of production; for the bond-addition experiment we equilibrate an $NVT$ liquid at $T = 80$~K and $\rho = 0.8\,\sigma^{-3}$ before inserting bonds and running for $5 \times 10^{5}$ production steps; for the non-interacting noise study we disable the pair potential entirely and run at low density and high temperature.
Harmonic bonds for the bond-addition experiment are added between random eligible pairs with stiffness $k_{b} = 100\,\varepsilon/\sigma^{2}$ and rest length set at the current pair separation.
When run-time bond insertion is required we use ESPResSo~\cite{TODO:espresso} in place of LAMMPS.

\paragraph{Water (Experiment~\ref{subsec:water-phases}).}
Water is modelled with TIP4P/2005~\cite{TODO:tip4p-2005}, with rigid geometry maintained by SHAKE.
We use $|\N_{\text{water}}| = 512$ molecules (1536 atoms) in a cubic box with periodic boundary conditions, a $10$~\AA{} real-space cutoff for the Lennard-Jones contribution, and a particle-particle particle-mesh Ewald treatment of long-range electrostatics with a real-space cutoff of $12$~\AA{}.
The integration timestep is $\delta t = 1$~fs.
The liquid branch is equilibrated in $NPT$ at $P = 1$~atm for $1$~ns and run for $5$~ns of production at each temperature in $\{240, 260, 280, 300, 320, 350, 400, 450, 500\}$~K; the ice branch at $240$~K is initialised from a proton-disordered ice Ih configuration and monitored through a pair correlation function to confirm the crystalline structure is retained.
Velocities are written every $20$~fs.

\paragraph{Molten salt (Experiment~\ref{subsec:molten-salt-ionic-conductivity}).}
Molten NaCl is modelled with the Tosi--Fumi pair potential~\cite{TODO:tosi-fumi} at $216$ ion pairs (432 atoms) in a cubic $NVT$ box set to the experimental density of $1.54$~g\,cm$^{-3}$.
Long-range electrostatics are again treated with PPPM.
The integration timestep is $\delta t = 1$~fs.
At each temperature in $\{1100, 1200, 1400, 1600, 1800\}$~K we equilibrate for $200$~ps and run for $2$~ns of production; velocities are written every $10$~fs.

\paragraph{Reproducibility.}
All input decks, analysis scripts, and random seeds used to generate the figures in Section~\ref{sec:experiments} are available in the repository described in Section~\ref{sec:data-and-availability}; each figure is regenerable from a single command and each run records its seed, source commit, and input hash in a provenance file.
